\documentclass[12pt]{article}
\usepackage[utf8]{inputenc}
\usepackage[T1]{fontenc}
\usepackage{hyperref}
\title{Talu Language Resources for K-12 Curriculum Development}
\author{Compiled by Claude Sonnet 5.5 (claude-sonnet-5-5) for Andr\'as Kornai}
\date{}
\begin{document}
\maketitle

\section{Introduction}
Talu (ISO 639-3 yta; Glottocode talu1238) is a Loloish (Lolo-Burmese) language of northern Yunnan, China. Wikipedia reports about 10,000 speakers in Yongsheng, Ninglang and Huaping counties, and a close relationship with Lolopo. One mission profile says that the language has never been specifically studied. Very little material was found, and nothing suggests it is taught in schools.

\section{Primary Educational Resources}
No primers, textbooks or school programmes were found.

\section{Linguistic and Reference Materials}
\begin{itemize}
\item Wikipedia, Talu language: \url{https://en.wikipedia.org/wiki/Talu_language}
\item Endangered Languages Project entry: \url{https://www.endangeredlanguages.com/lang/5172?hl=en}
\end{itemize}
Two further links appeared in the results, but their content was not examined and they may not concern Talu, so they are not relied on: \url{https://www.degruyterbrill.com/document/doi/10.1515/flih-2020-0006/html} and \url{https://www.academia.edu/5737342}.

\section{Religious and Community Resources}
\begin{itemize}
\item Asia Harvest people-group profile: \url{https://people-groups.asiaharvest.org/people-group-profiles/yunnan/talu.pdf}
\end{itemize}
No community-run language organization was found.

\section{Digital and Audio Resources}
No audio, video, corpus or software resources for Talu were found.

\section{Curriculum Development Context}
There are no existing teaching materials and no known writing system, so any curriculum would depend on new documentation and orthography work. The sources found here are descriptive only. The language is not known to be taught in schools.
\end{document}